\documentclass[11pt,a4paper]{article}

\usepackage[utf8]{inputenc}
\usepackage[T1]{fontenc}
\usepackage[french]{babel}
\usepackage{amsmath,amssymb}
\usepackage{graphicx}
\usepackage{booktabs}
\usepackage{hyperref}
\usepackage{geometry}
\usepackage{float}
\usepackage{xcolor}

\geometry{margin=2.5cm}
\definecolor{primary}{RGB}{46,134,171}

\title{
    \textbf{\color{primary}Prédiction de l'Âge Biologique par Méthylation de l'ADN} \\
    \large Analyse Computationnelle des Horloges Épigénétiques
}
\author{DNAm Age Prediction Benchmark}
\date{01 March 2026}

\begin{document}

\maketitle

\section*{Résumé}
Cette étude présente une analyse computationnelle approfondie de la prédiction de l'âge biologique à partir des profils de méthylation de l'ADN. Nous avons évalué 3 algorithmes d'apprentissage automatique sur une cohorte de 400 échantillons couvrant une plage d'âge de 18 à 90 ans.

Le modèle \textbf{Stack Optuna} a démontré les meilleures performances avec une erreur absolue moyenne (MAE) de 3.33 années et un coefficient de détermination (R\textsuperscript{2}) de 0.9309. La corrélation de Pearson entre l'âge prédit et l'âge chronologique atteint r = 0.9652 (p < 0.001), indiquant une très forte association linéaire.

L'analyse des résidus révèle une distribution approximativement normale (test de Shapiro-Wilk : W = 0.7417, p = 0.0000) avec un biais moyen de 0.27 années. Le coefficient de concordance de Lin (CCC = 0.9651) confirme l'excellente concordance entre les valeurs prédites et observées.

\tableofcontents
\newpage

\section{Introduction}
\subsection{Contexte scientifique}
La méthylation de l'ADN constitue l'une des modifications épigénétiques les plus étudiées, impliquant l'ajout covalent d'un groupe méthyle (CH\textsubscript{3}) sur le carbone 5 des cytosines au sein des dinucléotides CpG. Ce processus joue un rôle crucial dans la régulation de l'expression génique.

\subsection{Intérêt clinique}
L'écart entre l'âge épigénétique prédit et l'âge chronologique, appelé "accélération épigénétique" (EAA), représente un biomarqueur prometteur du vieillissement biologique. Une EAA positive a été associée à un risque accru de mortalité et de maladies liées à l'âge.

\section{Matériel et Méthodes}
L'échantillon comprend 400 individus. Le prétraitement a inclus le contrôle qualité, la normalisation des valeurs bêta, et l'imputation par KNN.

Six algorithmes ont été évalués : Elastic Net, Lasso, Ridge, Random Forest, XGBoost et AltumAge (MLP). La validation a été réalisée par partition 80/20 et validation croisée 5-fold.

\section{Résultats}
\subsection{Performance globale}

\begin{table}[H]
\centering
\begin{tabular}{lcccc}
\toprule
\textbf{Modèle} & \textbf{MAE} & \textbf{RMSE} & \textbf{R\textsuperscript{2}} & \textbf{Corrélation} \\
\midrule
Stack Optuna & 3.33 & 5.58 & 0.9309 & 0.9652 \\
Residual Learning & 3.33 & 5.89 & 0.9230 & 0.9614 \\
ElasticNetCV & 3.43 & 5.82 & 0.9247 & 0.9619 \\
\bottomrule
\end{tabular}
\caption{Performances des modèles de prédiction d'âge}
\end{table}

\subsection{Tests de comparaison}
Les tests t appariés confirment la supériorité statistique du modèle \textbf{Stack Optuna} :
\begin{itemize}
\item Stack Optuna vs Residual Learning : p = 0.9792
\item Stack Optuna vs ElasticNetCV : p = 0.2212
\end{itemize}

\subsection{Statistiques avancées}
\begin{table}[H]
\centering
\begin{tabular}{lll}
\toprule
\textbf{Métrique} & \textbf{Valeur} & \textbf{Interprétation} \\
\midrule
Corrélation de Pearson & 0.9652 & Très forte \\
CCC de Lin & 0.9651 & Excellente concordance \\
Biais moyen & 0.27 ans & Biais négligeable \\
Test de Shapiro-Wilk & p = 0.0000 & Non-normal \\
Durbin-Watson & 2.069 & Pas d'autocorrélation \\
\bottomrule
\end{tabular}
\end{table}

\section{Discussion}
Le modèle \textbf{Stack Optuna} offre une précision compatible avec les applications de recherche. La forte corrélation (r = 0.965) confirme la validité de notre approche.

L'accélération épigénétique constitue un biomarqueur précieux. Les individus présentant une EAA positive pourraient bénéficier d'interventions préventives ciblées.

\section{Conclusion}
Cette étude démontre l'efficacité des approches d'apprentissage automatique pour la quantification du vieillissement biologique. Les travaux futurs devront se concentrer sur la validation externe et l'intégration de données multi-omiques.

\end{document}
